\documentclass[11pt]{article}
\usepackage{ochamath}
\subjectcolor{2F5DA8}
\setsheet{線形代数・電気系院試補充}{リアプノフ方程式　演習}{https://university-math-crj.pages.dev/linear-algebra/lyapunov-equations/}
\namelinetrue
\begin{document}
\makesheethead{問題}
自作問題。本文の例題とは数値や設定が異なります。条件・途中式・検算を示してください。
\problem[連続時間の成分計算]
$A=\begin{pmatrix}-2&1\\0&-3\end{pmatrix},Q=I$ のリアプノフ方程式を解き、正定値を確認せよ。
\workspace{45mm}
\problem[積分表示]
安定な $A$ に対して $P=\int_0^\infty e^{A^{\mathsf T}t}Qe^{At}dt$ がリアプノフ方程式を満たすことを示せ。
\workspace{45mm}
\newpage
\problem[逆向きの安定性]
$P,Q\succ0$ と $A^*P+PA=-Q$ から全固有値の実部が負と証明せよ。
\workspace{45mm}
\problem[離散時間]
$A=\operatorname{diag}(1/3,1/2),Q=I$ に対し $A^{\mathsf T}PA-P=-I$ を解け。
\workspace{45mm}
\newpage
\problem[Sylvester方程式]
$A=\operatorname{diag}(1,2),B=\operatorname{diag}(3,4),C=\begin{pmatrix}4&10\\15&24\end{pmatrix}$ の $AX+XB=C$ を解け。
\workspace{45mm}
\problem[解はあっても正定値でない]
スカラー $A=2,Q=1$ で連続時間のリアプノフ方程式を解き、安定性との関係を述べよ。
\workspace{45mm}
\end{document}
